\documentclass[11pt]{article}
\usepackage[margin=1in]{geometry}
\usepackage{amsmath,amssymb}
\usepackage{hyperref}

\title{Literature Status: Hereditary Counts of Spreading Models}
\author{Run \texttt{fa\_banach\_001}}
\date{2026-06-23}

\begin{document}
\sloppy
\maketitle

\paragraph{Status.}
\textbf{Literature already answered.} This packet records that several
hereditary spreading-model questions from Androulakis--Odell--Schlumprecht--
Tomczak-Jaegermann, \emph{On the structure of the spreading models of a
Banach space}, arXiv:math/0305082, are explicitly answered in later
literature.

\paragraph{Source questions.}
The source questions are TeX labels \texttt{qst3.3} and \texttt{qst3.5} in
the arXiv source, printed as Questions 3.15 and 3.17 on pages 19--20 of the
source PDF. They ask, among other things:
\begin{itemize}
\item whether a Banach space whose every infinite-dimensional subspace admits
spreading models equivalent to the unit vector bases of \(\ell_1\) and
\(\ell_2\) must contain more, perhaps uncountably many, mutually
non-equivalent spreading models;
\item whether, for each \(n\in\mathbb N\), there is a Banach space whose every
subspace has exactly \(n\) different spreading models;
\item whether there is a Banach space whose every subspace has exactly
countably many different spreading models.
\end{itemize}

\paragraph{Supporting answer.}
Beanland--Freeman--Motakis, \emph{The stabilized set of \(p\)'s in Krivine's
theorem can be disconnected}, arXiv:1408.0265, gives the required examples.
Its Theorem 1, stated on page 2 of the supporting PDF, says that if
\[
F\subseteq [1,\infty]
\]
is either finite or an increasing sequence together with its limit, then there
is a reflexive Banach space \(X\) with an unconditional basis such that every
infinite-dimensional block subspace \(Y\) has the following behavior:
\begin{itemize}
\item \(\ell_p\) is finitely block represented in \(Y\) exactly for
\(p\in F\);
\item if \(F\) is finite, then the spreading models admitted by \(Y\) are
exactly the spaces \(\ell_p\), \(p\in F\);
\item if \(F\) is an increasing sequence with limit \(p_\omega\), then every
spreading model admitted by \(Y\) is isomorphic to \(\ell_p\) for some
\(p\in F\), and every non-limit \(p\in F\) is admitted.
\end{itemize}
On page 3, the same paper explicitly says that Theorem 1 solves the AOST
questions about exactly \(n\) many spreading models, exactly countably
infinitely many spreading models, and whether hereditary \(\ell_1\) and
\(\ell_2\) spreading models force uncountably many spreading models.

\paragraph{Consequences for the AOST questions.}
Taking \(F\) finite with \(|F|=n\) gives the exactly-\(n\) examples. Taking
\(F\) to be an increasing sequence together with its limit gives a countable
family of possible \(\ell_p\) spreading models, and the supporting authors
identify this as the solution to the countably-infinite question. Finally,
taking \(F=\{1,2\}\) gives a space whose hereditary subspace structure admits
\(\ell_1\) and \(\ell_2\) spreading models but not uncountably many different
ones, answering the uncountability alternative negatively.

\paragraph{Provenance.}
This is an explicit literature answer: Beanland--Freeman--Motakis cite AOST
and state that their Theorem 1 solves these open questions. The packet is
therefore not an original proof or a new partial result.

\paragraph{Scope limitations.}
This packet covers only the hereditary counting questions above. It does not
address AOST's other questions about stabilizing all spreading models, the
realisable partially ordered structures of \(SP_\omega(X)\), consequences of
finite or countable \(SP_\omega(X)\), or the unique-spreading-model question.

\begin{thebibliography}{9}

\bibitem{AOST}
G. Androulakis, E. Odell, Th. Schlumprecht, and N. Tomczak-Jaegermann,
\emph{On the structure of the spreading models of a Banach space},
arXiv:math/0305082; Canad. J. Math. 57 (2005), no. 4, 673--707.

\bibitem{BFM}
K. Beanland, D. Freeman, and P. Motakis,
\emph{The stabilized set of \(p\)'s in Krivine's theorem can be disconnected},
arXiv:1408.0265.

\end{thebibliography}

\end{document}
